\documentclass[11pt,aspectratio=169]{beamer}
\usetheme{Madrid}
\usecolortheme{default}

% Colores de la Facultad de Informática de la UMU (Granate)
\definecolor{informatica}{RGB}{121,11,21}
\setbeamercolor*{palette primary}{use=structure,fg=white,bg=informatica}
\setbeamercolor*{palette secondary}{use=structure,fg=white,bg=informatica!75!black}
\setbeamercolor*{palette tertiary}{use=structure,fg=white,bg=informatica!50!black}
\setbeamercolor{titlelike}{parent=palette primary}
\setbeamercolor{frametitle}{bg=informatica,fg=white}
\setbeamercolor{item}{fg=informatica}
\setbeamercolor{block title}{bg=informatica,fg=white}

\usepackage[utf8]{inputenc}
\usepackage[spanish]{babel}
\usepackage{graphicx}
\usepackage{listings}

% Información de la portada
\title[KGQA Chatbot]{Creación de un chatbot con IA capaz de buscar información en fuentes externas}
\subtitle{Trabajo Fin de Grado}
\author[M. Martínez Turpín]{Mario Martínez Turpín}
\institute[FIUM]{
  Facultad de Informática \\
  Universidad de Murcia
}
\date{Julio 2026} % Cambiar por la fecha real de la defensa

\begin{document}

% Diapositiva 1: Portada
\begin{frame}
  \titlepage
\end{frame}

% Diapositiva 2: Índice
\begin{frame}{Índice de la presentación}
  \tableofcontents
\end{frame}

% Diapositiva 3: Motivación
\section{Introducción}
\begin{frame}{Motivación del Proyecto}
  \begin{itemize}
    \item \textbf{El auge de los LLMs:} Modelos de lenguaje como ChatGPT son excelentes interfaces conversacionales, pero sufren de \textit{alucinaciones} cuando se les piden datos factuales y concretos.
    \vspace{0.3cm}
    \item \textbf{La rigidez de la Web Semántica:} Wikidata y las bases RDF son extremadamente precisas, pero requieren conocer lenguajes de consulta complejos (SPARQL).
    \vspace{0.3cm}
    \item \textbf{La Solución (KGQA):} Crear un sistema de pregunta-respuesta sobre Grafos de Conocimiento que una la accesibilidad del lenguaje natural con la precisión matemática del esquema semántico.
  \end{itemize}
\end{frame}

% Diapositiva 4: Objetivos
\section{Objetivos}
\begin{frame}{Objetivos del TFG}
  \begin{enumerate}
    \item Desarrollar una interfaz conversacional intuitiva y transparente.
    \vspace{0.2cm}
    \item Traducir consultas de texto libre a código \textbf{SPARQL} válido limitando las propiedades a las permitidas en la ontología.
    \vspace{0.2cm}
    \item Construir un \textbf{Grafo de Conocimiento local} del dominio del cine para garantizar rapidez y privacidad frente a APIs de terceros.
    \vspace{0.2cm}
    \item Implementar un \textbf{Agente Autónomo (LangChain)} capaz de detectar y autocorregir sus propios errores de ejecución en tiempo real.
  \end{enumerate}
\end{frame}

% Diapositiva 5: Arquitectura
\section{Arquitectura del Sistema}
\begin{frame}{Arquitectura en 3 Capas}
  \begin{columns}
    \begin{column}{0.5\textwidth}
      La arquitectura \textbf{RAG (Retrieval-Augmented Generation)} se divide en:
      \begin{itemize}
        \item \textbf{Frontend (Streamlit):} Interfaz gráfica moderna, historial de chat y panel desplegable de explicabilidad (Explainable AI).
        \item \textbf{Lógica y Agente (LangChain):} Cerebro del sistema, extracción de entidades y orquestación del LLM.
        \item \textbf{Datos (RDFLib):} Grafo local en \texttt{.ttl} procesado a partir de Wikidata.
      \end{itemize}
    \end{column}
    \begin{column}{0.5\textwidth}
      % TODO: Incluir la imagen del diagrama de arquitectura aquí cuando la hagas
      \begin{center}
        \includegraphics[width=0.9\textwidth, height=4cm, keepaspectratio]{example-image} \\
        \textit{\small *Reemplazar por el diagrama de arquitectura*}
      \end{center}
    \end{column}
  \end{columns}
\end{frame}

% Diapositiva 6: Flujo del Agente
\section{Bucle de Reflexión del Agente}
\begin{frame}{El Bucle de Reflexión Autocorrector}
  \begin{block}{El Problema}
    La sintaxis SPARQL es estricta. Un LLM suele cometer errores de puntuación que provocan fallos catastróficos.
  \end{block}
  
  \begin{block}{La Solución: Flujo Agéntico}
    \begin{enumerate}
      \item El Agente genera una consulta SPARQL inicial.
      \item El motor RDFLib intenta ejecutarla localmente.
      \item Si hay un error, el sistema \textbf{no aborta}. Atrapa la excepción.
      \item Se devuelve el mensaje de error de compilación al Agente.
      \item El Agente reflexiona y regenera la consulta ya corregida de forma invisible para el usuario.
    \end{enumerate}
  \end{block}
\end{frame}

% Diapositiva 7: Evaluación y Resultados
\section{Evaluación y Resultados}
\begin{frame}{Evaluación y Resultados}
  Se evaluó el sistema automáticamente con un \textit{dataset} de 20 preguntas complejas (directores, años de publicación, géneros, etc.).
  
  \vspace{0.4cm}
  \begin{table}[]
    \centering
    \begin{tabular}{|l|c|c|}
    \hline
    \textbf{Métrica Evaluada} & \textbf{Aciertos} & \textbf{Precisión (\%)} \\ \hline
    Extracción y Desambiguación de Entidad & 16 / 20 & \textbf{80.00\%} \\ \hline
    Generación de SPARQL Válido & 15 / 20 & \textbf{75.00\%} \\ \hline
    \end{tabular}
  \end{table}
  \vspace{0.3cm}

  \textbf{Discusión:}
  \begin{itemize}
    \item El fallo del 20\% en entidades se debe a ambigüedades idiomáticas (Ej: "Origen" vs "Inception") y los límites de peticiones HTTP.
    \item Un \textbf{75\% de éxito} en SPARQL demuestra que inyectar el diccionario de propiedades en el \textit{prompt} previene alucinaciones masivas.
  \end{itemize}
\end{frame}

% Diapositiva 8: Conclusiones
\section{Conclusiones}
\begin{frame}{Conclusiones y Trabajo Futuro}
  \textbf{Conclusiones Principales:}
  \begin{itemize}
    \item El uso de Agentes LangChain soluciona las barreras de uso de la Web Semántica para usuarios sin perfil técnico.
    \item El bucle de reflexión eleva drásticamente la tolerancia a fallos del sistema sin comprometer la latencia.
  \end{itemize}
  
  \vspace{0.5cm}
  \textbf{Trabajo Futuro:}
  \begin{itemize}
    \item Desambiguación mediante búsqueda vectorial (\textit{embeddings}) fonética.
    \item Integración de \textit{Few-Shot Prompting} dinámico para soportar consultas \textit{Multi-Salto} complejas.
  \end{itemize}
\end{frame}

% Diapositiva 9: Fin
\begin{frame}
  \begin{center}
    \Huge \textbf{Gracias por su atención}\\
    \vspace{1.5cm}
    \Large ¿Alguna pregunta?
  \end{center}
\end{frame}

\end{document}
